The quasilinear diffusion operator for radiofrequency was developped by
Kennel and Engelmann. The relativistic treatment was performed by Lerche,
who proposes the following operator 
\begin{eqnarray}
Q\left( f\right)  &=&-\frac{e^{2}}{\left( 2\pi \right) ^{3}}%
\lim_{V\rightarrow \infty }\frac{1}{V}\sum_{n=-\infty }^{\infty }\iiint
d^{3}k\left[ \left( P_{\Vert }^{*}E_{\mathbf{k},\Vert }^{*}J_{n}+\left[
P_{\bot }^{*}-\frac{1}{p_{\bot }}\left( \frac{k_{\Vert }v_{\Vert }}{\omega
^{*}}-1\right) \right] E_{\mathbf{k},\bot }^{*}\right) \right. 
\label{qloporig} \\
&&\left. \cdot \frac{i}{\left[ n\Omega +k_{\Vert }v_{\Vert }-\omega \right] }%
\left( E_{\mathbf{k},\Vert }J_{n}P_{\Vert }+E_{\mathbf{k},\bot }P_{\bot
}\right) f\right]   \nonumber
\end{eqnarray}
with 
\begin{eqnarray}
P_{\Vert } &=&\frac{\partial }{\partial p_{\Vert }}-\frac{n\Omega }{\omega
v_{\bot }}\left( v_{\bot }\frac{\partial }{\partial p_{\Vert }}-v_{\Vert }%
\frac{\partial }{\partial p_{\bot }}\right)   \label{ppar} \\
P_{\bot } &=&\frac{\partial }{\partial p_{\bot }}+\frac{k_{\Vert }}{\omega }%
\left( v_{\bot }\frac{\partial }{\partial p_{\Vert }}-v_{\Vert }\frac{%
\partial }{\partial p_{\bot }}\right)   \label{pperp} \\
E_{\mathbf{k},\bot } &=&\frac{1}{\sqrt{2}}\left( E_{\mathbf{k},-}e^{i\alpha
}J_{n+1}+E_{\mathbf{k},+}e^{-i\alpha }J_{n-1}\right)   \label{eperp}
\end{eqnarray}

The electric field is decomposed into its Fourier components 
\begin{equation}
\mathbf{E}(\mathbf{x},t)=\iiint \frac{d^{3}k}{(2\pi )^{3}}e^{i\mathbf{k}%
\cdot \mathbf{x}-i\omega t}\mathbf{E}_{\mathbf{k}}\left( \mathbf{k}\right)
\end{equation}
which are projected on the rotating field frame 
\begin{eqnarray}
E_{\mathbf{k},\pm } &=&\frac{1}{\sqrt{2}}\left( E_{\mathbf{k},x}\pm iE_{%
\mathbf{k},y}\right)  \label{ekpm} \\
E_{\mathbf{k},\Vert } &=&E_{\mathbf{k},z}
\end{eqnarray}

The wave vector is expressed in cylindrical coordinates as 
\begin{eqnarray}
k_{x} &=&k_{\bot }\cos \alpha \\
k_{y} &=&k_{\bot }\sin \alpha  \nonumber \\
k_{z} &=&k_{\Vert }  \nonumber
\end{eqnarray}
and the argument of the Bessel functions is $k_{\bot }v_{\bot }/\Omega $.
The relativistic cyclotron frequency is 
\begin{equation}
\Omega =\frac{q_{e}B}{\gamma m_{e}}
\end{equation}

\subsection{Conservative formulation of the RF wave operator}

In order to incorporate the RF wave physics in the Fokker-Planck or Drift
Kinetic solvers, the operator (\ref{qloporig}) must be cast in a
conservative form.

\subsubsection{Spherical coordinates representation}

The transformation from cylindrical to spherical momentum derivatives is
given by 
\begin{eqnarray}
\frac{\partial }{\partial p_{\Vert }} &=&\frac{\partial p}{\partial p_{\Vert
}}\frac{\partial }{\partial p}+\frac{\partial \xi }{\partial p_{\Vert }}%
\frac{\partial }{\partial \xi }  \label{dppar} \\
\frac{\partial }{\partial p_{\bot }} &=&\frac{\partial p}{\partial p_{\bot }}%
\frac{\partial }{\partial p}+\frac{\partial \xi }{\partial p_{\bot }}\frac{%
\partial }{\partial \xi }  \label{dpperp}
\end{eqnarray}

Since 
\begin{eqnarray}
p &=&\sqrt{p_{\Vert }^{2}+p_{\bot }^{2}} \\
\xi &=&\frac{p_{\Vert }}{\sqrt{p_{\Vert }^{2}+p_{\bot }^{2}}}
\end{eqnarray}
one obtains 
\begin{eqnarray}
\frac{\partial p}{\partial p_{\Vert }} &=&\xi \\
\frac{\partial p}{\partial p_{\bot }} &=&\sqrt{1-\xi ^{2}} \\
\frac{\partial \xi }{\partial p_{\Vert }} &=&\frac{1-\xi ^{2}}{p} \\
\frac{\partial \xi }{\partial p_{\bot }} &=&-\frac{\xi \sqrt{1-\xi ^{2}}}{p}
\end{eqnarray}

Hence rewritting (\ref{ppar}-\ref{pperp}) in cylindrical coordinates using (%
\ref{dppar}-\ref{dpperp}) 
\begin{eqnarray}
P_{\Vert } &=&\xi \frac{\partial }{\partial p}+\frac{1}{p}\left( 1-\xi ^{2}-%
\frac{n\Omega }{\omega }\right) \frac{\partial }{\partial \xi }
\label{ppar2} \\
P_{\bot } &=&\sqrt{1-\xi ^{2}}\frac{\partial }{\partial p}-\frac{\sqrt{1-\xi
^{2}}}{p\xi }\left( \xi ^{2}-\frac{k_{\Vert }v_{\Vert }}{\omega }\right) 
\frac{\partial }{\partial \xi }  \label{pperp2}
\end{eqnarray}

In order to permute derivatives and the integral in expression (\ref
{qloporig}) and obtain a conservative formulation, one must express both $%
P_{\Vert }$ and $P_{\bot }$ in terms of the derivatives $\frac{1}{p^{2}}%
\frac{\partial }{\partial p}p^{2}$ and $-\frac{1}{p}\frac{\partial }{%
\partial \xi }.$ Since 
\begin{eqnarray}
\xi \frac{\partial }{\partial p} &=&\frac{1}{p^{2}}\frac{\partial }{\partial
p}p^{2}\xi -\frac{2\xi }{p} \\
\sqrt{1-\xi ^{2}}\frac{\partial }{\partial p} &=&\frac{1}{p^{2}}\frac{%
\partial }{\partial p}p^{2}\sqrt{1-\xi ^{2}}-\frac{2\sqrt{1-\xi ^{2}}}{p}
\end{eqnarray}
and 
\begin{eqnarray}
\frac{1}{p}\left( 1-\xi ^{2}-\frac{n\Omega }{\omega }\right) \frac{\partial 
}{\partial \xi } &=&-\frac{1}{p}\frac{\partial }{\partial \xi }\left( -1+\xi
^{2}+\frac{n\Omega }{\omega }\right) +\frac{2\xi }{p} \\
-\frac{\sqrt{1-\xi ^{2}}}{p\xi }\left( \xi ^{2}-\frac{k_{\Vert }v_{\Vert }}{%
\omega }\right) \frac{\partial }{\partial \xi } &=&-\frac{1}{p}\frac{%
\partial }{\partial \xi }\frac{\sqrt{1-\xi ^{2}}}{\xi }\left( \xi ^{2}-\frac{%
k_{\Vert }v_{\Vert }}{\omega }\right)  \nonumber \\
&&+\frac{1}{p_{\bot }}\left( \frac{k_{\Vert }v_{\Vert }}{\omega }-1\right) +%
\frac{2\sqrt{1-\xi ^{2}}}{p}
\end{eqnarray}
the operators in (\ref{qloporig}) can be expressed in a divergence form 
\begin{eqnarray}
P_{\Vert } &=&\frac{1}{p^{2}}\frac{\partial }{\partial p}p^{2}\xi -\frac{1}{p%
}\frac{\partial }{\partial \xi }\left( -1+\xi ^{2}+\frac{n\Omega }{\omega }%
\right) \\
P_{\bot }-\frac{1}{p_{\bot }}\left( \frac{k_{\Vert }v_{\Vert }}{\omega }%
-1\right) &=&\frac{1}{p^{2}}\frac{\partial }{\partial p}p^{2}\sqrt{1-\xi ^{2}%
}-\frac{1}{p}\frac{\partial }{\partial \xi }\frac{\sqrt{1-\xi ^{2}}}{\xi }%
\left( \xi ^{2}-\frac{k_{\Vert }v_{\Vert }}{\omega }\right)
\end{eqnarray}
so that (\ref{qloporig}) can be rewritten as 
\begin{equation}
Q\left( f\right) =-\nabla \cdot S^{\text{RF}}=-\frac{1}{p^{2}}\frac{\partial 
}{\partial p}\left( p^{2}S_{p}^{\text{RF}}\right) +\frac{1}{p}\frac{\partial 
}{\partial \xi }\left( \sqrt{1-\xi ^{2}}S_{\xi }^{\text{RF}}\right)
\end{equation}
with 
\begin{eqnarray}
S_{p}^{\text{RF}} &=&\lim_{V\rightarrow \infty }\frac{e^{2}}{\left( 2\pi
\right) ^{3}}\frac{1}{V}\sum_{n=-\infty }^{\infty }\iiint d^{3}k\left[
\left( \xi E_{\mathbf{k},\Vert }^{*}J_{n}+\sqrt{1-\xi ^{2}}E_{\mathbf{k}%
,\bot }^{*}\right) \right.  \nonumber \\
&&\left. \cdot \frac{i}{\left[ n\Omega +k_{\Vert }v_{\Vert }-\omega \right] }%
\left( E_{\mathbf{k},\Vert }J_{n}P_{\Vert }+E_{\mathbf{k},\bot }P_{\bot
}\right) f\right]  \label{fluxRFp}
\end{eqnarray}
\begin{eqnarray}
S_{\xi }^{\text{RF}} &=&\lim_{V\rightarrow \infty }\frac{e^{2}}{\left( 2\pi
\right) ^{3}}\frac{1}{V}\sum_{n=-\infty }^{\infty }\iiint d^{3}k\left[
\left( \frac{-1}{\sqrt{1-\xi ^{2}}}\left( 1-\xi ^{2}-\frac{n\Omega }{\omega
^{*}}\right) E_{\mathbf{k},\Vert }^{*}J_{n}+\frac{1}{\xi }\left( \xi ^{2}-%
\frac{k_{\Vert }v_{\Vert }}{\omega ^{*}}\right) E_{\mathbf{k},\bot
}^{*}\right) \right.  \nonumber \\
&&\left. \cdot \frac{i}{\left[ n\Omega +k_{\Vert }v_{\Vert }-\omega \right] }%
\left( E_{\mathbf{k},\Vert }J_{n}P_{\Vert }+E_{\mathbf{k},\bot }P_{\bot
}\right) f\right]  \label{fluxRFksi}
\end{eqnarray}

Considering the above expressions of the fluxes and the operators (\ref
{ppar2}-\ref{pperp2}), momentum space dynamics has only a diffusive nature,
and therefore 
\begin{equation}
\mathbf{S}^{\text{RF}}=-\Bbb{D}^{\text{RF}}\cdot \nabla _{p}f
\end{equation}
with 
\begin{equation}
\Bbb{D}^{\text{RF}}=\left( 
\begin{array}{ll}
D_{pp}^{\text{RF}} & D_{p\xi }^{\text{RF}} \\ 
D_{\xi p}^{\text{RF}} & D_{\xi \xi }^{\text{RF}}
\end{array}
\right)
\end{equation}

Hence, using the expression of the gradient in spherical coordinates, 
\begin{eqnarray}
S_{p}^{\text{RF}} &=&-D_{pp}^{\text{RF}}\frac{\partial f}{\partial p}%
+D_{p\xi }^{\text{RF}}\frac{\sqrt{1-\xi ^{2}}}{p}\frac{\partial f}{\partial
\xi } \\
S_{\xi }^{\text{RF}} &=&-D_{\xi p}^{\text{RF}}\frac{\partial f}{\partial p}%
+D_{\xi \xi }^{\text{RF}}\frac{\sqrt{1-\xi ^{2}}}{p}\frac{\partial f}{%
\partial \xi }
\end{eqnarray}
with 
\begin{eqnarray}
D_{pp}^{\text{RF}} &=&-\lim_{V\rightarrow \infty }\frac{e^{2}}{\left( 2\pi
\right) ^{3}}\frac{1}{V}\sum_{n=-\infty }^{\infty }\iiint d^{3}k\left[
\left( \xi E_{\mathbf{k},\Vert }^{*}J_{n}+\sqrt{1-\xi ^{2}}E_{\mathbf{k}%
,\bot }^{*}\right) \right. \\
&&\left. \cdot \frac{i}{\left[ n\Omega +k_{\Vert }v_{\Vert }-\omega \right] }%
\left( \xi E_{\mathbf{k},\Vert }J_{n}+\sqrt{1-\xi ^{2}}E_{\mathbf{k},\bot
}\right) \right]  \nonumber \\
D_{p\xi }^{\text{RF}} &=&\lim_{V\rightarrow \infty }\frac{e^{2}}{\left( 2\pi
\right) ^{3}}\frac{1}{V}\sum_{n=-\infty }^{\infty }\iiint d^{3}k\left[
\left( \xi E_{\mathbf{k},\Vert }^{*}J_{n}+\sqrt{1-\xi ^{2}}E_{\mathbf{k}%
,\bot }^{*}\right) \right. \\
&&\left. \cdot \frac{i}{\left[ n\Omega +k_{\Vert }v_{\Vert }-\omega \right] }%
\left( \frac{1}{\sqrt{1-\xi ^{2}}}\left( 1-\xi ^{2}-\frac{n\Omega }{\omega }%
\right) E_{\mathbf{k},\Vert }J_{n}-\frac{1}{\xi }\left( \xi ^{2}-\frac{%
k_{\Vert }v_{\Vert }}{\omega }\right) E_{\mathbf{k},\bot }\right) \right] 
\nonumber \\
D_{\xi p}^{\text{RF}} &=&-\lim_{V\rightarrow \infty }\frac{e^{2}}{\left(
2\pi \right) ^{3}}\frac{1}{V}\sum_{n=-\infty }^{\infty }\iiint d^{3}k\left[
\left( \frac{-1}{\sqrt{1-\xi ^{2}}}\left( 1-\xi ^{2}-\frac{n\Omega }{\omega
^{*}}\right) E_{\mathbf{k},\Vert }^{*}J_{n}+\frac{1}{\xi }\left( \xi ^{2}-%
\frac{k_{\Vert }v_{\Vert }}{\omega ^{*}}\right) E_{\mathbf{k},\bot
}^{*}\right) \right. \\
&&\left. \cdot \frac{i}{\left[ n\Omega +k_{\Vert }v_{\Vert }-\omega \right] }%
\left( E_{\mathbf{k},\Vert }J_{n}\xi +E_{\mathbf{k},\bot }\sqrt{1-\xi ^{2}}%
\right) \right]  \nonumber \\
D_{\xi \xi }^{\text{RF}} &=&\lim_{V\rightarrow \infty }\frac{e^{2}}{\left(
2\pi \right) ^{3}}\frac{1}{V}\sum_{n=-\infty }^{\infty }\iiint d^{3}k\left[
\left( \frac{-1}{\sqrt{1-\xi ^{2}}}\left( 1-\xi ^{2}-\frac{n\Omega }{\omega
^{*}}\right) E_{\mathbf{k},\Vert }^{*}J_{n}+\frac{1}{\xi }\left( \xi ^{2}-%
\frac{k_{\Vert }v_{\Vert }}{\omega ^{*}}\right) E_{\mathbf{k},\bot
}^{*}\right) \right. \\
&&\left. \cdot \frac{i}{\left[ n\Omega +k_{\Vert }v_{\Vert }-\omega \right] }%
\left( \frac{1}{\sqrt{1-\xi ^{2}}}\left( 1-\xi ^{2}-\frac{n\Omega }{\omega }%
\right) E_{\mathbf{k},\Vert }J_{n}-\frac{1}{\xi }\left( \xi ^{2}-\frac{%
k_{\Vert }v_{\Vert }}{\omega }\right) E_{\mathbf{k},\bot }\right) \right] 
\nonumber
\end{eqnarray}

\bigskip In the limit of a resonant diffusion, 
\begin{equation}
\frac{1}{\left[ n\Omega +k_{\Vert }v_{\Vert }-\omega \right] }\rightarrow
i\pi \delta \left( \omega -k_{\Vert }v_{\Vert }-n\Omega \right)
\end{equation}
and with the resonance condition 
\begin{equation}
k_{\Vert }v_{\Vert }=\omega -n\Omega
\end{equation}
RF diffusion coefficients may be expressed in a simple form 
\begin{eqnarray}
D_{pp}^{\text{RF}} &=&\sum_{n=-\infty }^{\infty }\left( 1-\xi ^{2}\right)
D_{n}^{\text{RF}}\left( p,\xi \right) \\
D_{p\xi }^{\text{RF}} &=&\sum_{n=-\infty }^{\infty }-\frac{\sqrt{1-\xi ^{2}}%
}{\xi }\left( 1-\xi ^{2}-\frac{n\Omega }{\omega }\right) D_{n}^{\text{RF}%
}\left( p,\xi \right) \\
D_{\xi p}^{\text{RF}} &=&\sum_{n=-\infty }^{\infty }-\frac{\sqrt{1-\xi ^{2}}%
}{\xi }\left( 1-\xi ^{2}-\frac{n\Omega }{\omega }\right) D_{n}^{\text{RF}%
}\left( p,\xi \right) \\
D_{\xi \xi }^{\text{RF}} &=&\sum_{n=-\infty }^{\infty }\frac{1}{\xi ^{2}}%
\left( 1-\xi ^{2}-\frac{n\Omega }{\omega }\right) ^{2}D_{n}^{\text{RF}%
}\left( p,\xi \right)
\end{eqnarray}
with 
\begin{equation}
D_{n}^{\text{RF}}\left( p,\xi \right) =\lim_{V\rightarrow \infty }\frac{\pi
e^{2}}{V}\iiint \frac{d^{3}k}{\left( 2\pi \right) ^{3}}\delta \left( \omega
-k_{\Vert }v_{\Vert }-n\Omega \right) \left| E_{\mathbf{k},\bot }+\frac{\xi 
}{\sqrt{1-\xi ^{2}}}E_{\mathbf{k},\Vert }J_{n}\right| ^{2}  \label{dnrf}
\end{equation}

Using (\ref{eperp}), we can define 
\begin{equation}
D_{\mathbf{k}}^{(n)}=e^{2}\left| \frac{1}{\sqrt{2}}E_{\mathbf{k}%
,+}e^{-i\alpha }J_{n-1}\left( \frac{k_{\bot }v_{\bot }}{\Omega }\right) +%
\frac{1}{\sqrt{2}}E_{\mathbf{k},-}e^{+i\alpha }J_{n+1}\left( \frac{k_{\bot
}v_{\bot }}{\Omega }\right) +\frac{p_{\Vert }}{p_{\bot }}E_{\mathbf{k},\Vert
}J_{n}\left( \frac{k_{\bot }v_{\bot }}{\Omega }\right) \right| ^{2}
\label{dk}
\end{equation}
so that (\ref{dnrf}) becomes 
\begin{equation}
D_{n}^{\text{RF}}(p,\xi )=\lim_{V\rightarrow \infty }\frac{\pi }{V}\int 
\frac{d^{3}k}{(2\pi )^{3}}D_{\mathbf{k}}^{(n)}\delta (\omega -k_{\Vert
}v_{\Vert }-n\Omega )
\end{equation}

If rays are gathered in RF beams of different frequencies $\omega _{b}$,
diffusion coefficients are sums of each contribution ower all hamonics $n$ ,
leading to the expression

\begin{eqnarray}
D_{pp}^{\text{RF}} &=&\sum_{n=-\infty }^{+\infty }\sum_{b}(1-\xi
^{2})D_{b,n}^{\text{RF}}(p,\xi )  \label{dtensor1} \\
D_{p\xi }^{\text{RF}} &=&\sum_{n=-\infty }^{+\infty }\sum_{b}-\frac{\sqrt{%
1-\xi ^{2}}}{\xi }\left[ 1-\xi ^{2}-\frac{n\Omega }{\omega _{b}}\right]
D_{b,n}^{\text{RF}}(p,\xi )  \label{dtensor2} \\
D_{\xi p}^{\text{RF}} &=&\sum_{n=-\infty }^{+\infty }\sum_{b}-\frac{\sqrt{%
1-\xi ^{2}}}{\xi }\left[ 1-\xi ^{2}-\frac{n\Omega }{\omega _{b}}\right]
D_{b,n}^{\text{RF}}(p,\xi )  \label{dtensor3} \\
D_{\xi \xi }^{\text{RF}} &=&\sum_{n=-\infty }^{+\infty }\sum_{b}\frac{1}{\xi
^{2}}\left[ 1-\xi ^{2}-\frac{n\Omega }{\omega _{b}}\right] ^{2}D_{b,n}^{%
\text{RF}}(p,\xi )  \label{dtensor4}
\end{eqnarray}
where 
\begin{equation}
D_{n,b}^{\text{RF}}(p,\xi )=\lim_{V\rightarrow \infty }\frac{\pi }{V}\int 
\frac{d^{3}k}{(2\pi )^{3}}D_{\mathbf{k}}^{b,(n)}\delta (\omega _{b}-k_{\Vert
}v_{\Vert }-n\Omega )
\end{equation}
and $D_{\mathbf{k}}^{b,(n)}$ accounts for the polarization and the intensity
of the RF wave. It is given by 
\begin{equation}
D_{\mathbf{k}}^{b,(n)}=e^{2}\left| \frac{1}{\sqrt{2}}E_{\mathbf{k}
,b,+}e^{-i\alpha }J_{n-1}\left( \frac{k_{\bot ,b}v_{\bot }}{\Omega }\right) +%
\frac{1}{\sqrt{2}}E_{\mathbf{k},b,-}e^{+i\alpha }J_{n+1}\left( \frac{k_{\bot
,b}v_{\bot }}{\Omega }\right) +\frac{p_{\Vert }}{p_{\bot }}E_{\mathbf{k}%
,b,\Vert }J_{n}\left( \frac{k_{\bot ,b}v_{\bot }}{\Omega }\right) \right|
^{2}
\end{equation}

\subsubsection{Cylindrical coordinates representation}

The Fokker-Planck equation is usually solved numerically in spherical
coordinates, because of the spherical symmetry of the collisional operator.
Therefore, in this coordinate system, the numerical problem is
well-conditionned, ensuring stable convergence towards the solution.
However, in many case the RF quasilinear operator has a more cylindrical
symmetry, and it could be useful to derive its expression in this coordinate
system, in order to get a more physical insight of the wave-particle
interaction. Starting from the general expression of the flux divergence in
cylindrical coordinates, 
\begin{equation}
\nabla _{p}\cdot \mathbf{S}=\frac{1}{p_{\bot }}\frac{\partial }{\partial
p_{\bot }}\left( p_{\bot }S_{\bot }\right) +\frac{\partial }{\partial
p_{\Vert }}\left( S_{\Vert }\right)  \label{divcyl}
\end{equation}
and taking into account of the diffusive nature of the wave-particle
interaction, 
\begin{equation}
\mathbf{S}^{\text{RF}}=-\Bbb{D}^{\text{RF}}\cdot \nabla f=-\left( 
\begin{array}{cc}
D_{\bot \bot } & D_{\bot \Vert } \\ 
D_{\Vert \bot } & D_{\Vert \Vert }
\end{array}
\right) \cdot \left( 
\begin{array}{c}
\partial f/\partial p_{\bot } \\ 
\partial f/\partial p_{\Vert }
\end{array}
\right)
\end{equation}
where the cylindrical tensor elements are related to the spherical ones by 
\begin{equation}
\left( 
\begin{array}{c}
D_{\bot \bot } \\ 
D_{\bot \Vert } \\ 
D_{\Vert \bot } \\ 
D_{\Vert \Vert }
\end{array}
\right) =\left( 
\begin{array}{cccc}
\left( 1-\xi ^{2}\right) & \xi \sqrt{1-\xi ^{2}} & \xi \sqrt{1-\xi ^{2}} & 
\xi ^{2} \\ 
\xi \sqrt{1-\xi ^{2}} & -\left( 1-\xi ^{2}\right) & \xi ^{2} & -\xi \sqrt{%
1-\xi ^{2}} \\ 
\xi \sqrt{1-\xi ^{2}} & \xi ^{2} & -\left( 1-\xi ^{2}\right) & -\xi \sqrt{%
1-\xi ^{2}} \\ 
\xi ^{2} & -\xi \sqrt{1-\xi ^{2}} & -\xi \sqrt{1-\xi ^{2}} & \left( 1-\xi
^{2}\right)
\end{array}
\right) \cdot \left( 
\begin{array}{c}
D_{pp} \\ 
D_{p\xi } \\ 
D_{\xi p} \\ 
D_{\xi \xi }
\end{array}
\right)
\end{equation}

Applying this transformation, one obtains, 
\begin{eqnarray}
D_{\bot \bot }^{\text{RF}} &=&\sum_{n=-\infty }^{+\infty }\sum_{b}\left[ 
\frac{n\Omega }{\omega _{b}}\right] ^{2}D_{b,n}^{\text{RF}}(p,\xi )
\label{dcyl} \\
D_{\bot \Vert }^{\text{RF}} &=&\sum_{n=-\infty }^{+\infty }\sum_{b}\frac{%
\sqrt{1-\xi ^{2}}}{\xi }\frac{n\Omega }{\omega _{b}}\left[ 1-\frac{n\Omega }{%
\omega _{b}}\right] D_{b,n}^{\text{RF}}(p,\xi ) \\
D_{\Vert \bot }^{\text{RF}} &=&\sum_{n=-\infty }^{+\infty }\sum_{b}\frac{%
\sqrt{1-\xi ^{2}}}{\xi }\frac{n\Omega }{\omega _{b}}\left[ 1-\frac{n\Omega }{%
\omega _{b}}\right] D_{b,n}^{\text{RF}}(p,\xi ) \\
D_{\Vert \Vert }^{\text{RF}} &=&\sum_{n=-\infty }^{+\infty }\sum_{b}\frac{%
(1-\xi ^{2})}{\xi ^{2}}\left[ 1-\frac{n\Omega }{\omega _{b}}\right]
^{2}D_{b,n}^{\text{RF}}(p,\xi )
\end{eqnarray}

For $n=0$ (which corresponds to the Lower Hybrid wave) the quasilinear
diffusion is strictly along the parallel direction (i.e. magnetic field
line). Also, at a cyclotron harmonic, where $\omega _{b}=n\Omega $, the
diffusion is only perpendicular, when relativistic corrections are not
considered.
